\section{Symmetric replicas and entropy metrics}\label{sec:replicas}

We now lift the one-copy estimates of Section~\ref{sec:conditional} to
positive operators on symmetric tensor powers.  There are two outputs.
A transfer of a subsystem gives an operator lower bound with its
conditional mutual information in the exponent.  A similarity transform
of a local term has a scalar symbol whose departure from the nonnegative
real axis is controlled by Lemma~\ref{lem:skew}.  The latter assertion
requires a product calculus:
it is not enough to compute the expectation of the similarity transform
alone.

Throughout this section, $\mathcal V$ is a fixed finite-dimensional
complex Hilbert space with a specified tensor factorization, possibly
including auxiliary factors.  Write
\[
 d=\dim\mathcal V,\qquad
 \mathcal S_k=\Sym^k\mathcal V,\qquad
 D_k=\dim\mathcal S_k=\binom{k+d-1}{d-1},
\]
and let $\Pi_k$ be the orthogonal projector onto $\mathcal S_k$.
A subsystem $Q$ is a union of the specified factors; its complement is
$Q^c$, and $d_Q$ denotes its Hilbert-space dimension.  The empty
subsystem has dimension one.  An operator on $Q^{\otimes k}$ is extended
by the identity on the other factors.  The permutation of the $k$
copies of $Q$ associated with $\pi\in S_k$ is denoted by $U_Q(\pi)$;
$U(\pi)=U_{\mathcal V}(\pi)$ permutes entire copies.

Every limit below is taken with $\mathcal V$, its factorization, and all
real exponents fixed.  When these results are applied at a scale $n$,
the entire finite system, auxiliary dimensions, and finite collection of
partitions are fixed before $k\to\infty$.  Constants and degrees in
$\poly(k)$ may depend on these fixed data.  Errors $o_k(1)$ are uniform
over all density matrices supported on $\mathcal S_k$.  In particular,
they remain uniform for density matrices depending on an interpolation
parameter.  A cutoff and a polynomial approximation are fixed before
taking $k\to\infty$; their errors are removed only afterwards.

\subsection{Schur labels and their quantitative bounds}

Young-diagram labels encode the spectrum of an independent tensor-power
state in the spectrum-estimation theorem of Keyl and
Werner~\cite{KeylWerner2001}.  Here we need label estimates for arbitrary
permutation-invariant densities as well as tensor powers.  We derive the
specific inequalities from Schur--Weyl duality below.

A partition $\lambda\vdash k$ is a decreasing sequence of nonnegative
integers summing to $k$.  Trailing zeroes may be appended.  Let
$[\lambda]$ be the corresponding irreducible representation of $S_k$,
and write $d_\lambda=\dim[\lambda]$.  Schur--Weyl duality and the Weyl
character formula give
\begin{equation}\label{replicas:schur-decomposition}
 (\C^q)^{\otimes k}
 =\bigoplus_{\substack{\lambda\vdash k\\\ell(\lambda)\le q}}
       [\lambda]\otimes V^{(q)}_\lambda,
 \qquad
 \chi^{(q)}_\lambda(x)
 =\frac{\det[x_i^{\lambda_j+q-j}]}{\det[x_i^{q-j}]}.
\end{equation}
Here the two factors are the irreducible modules for the permutation
and diagonal general-linear actions, respectively, and the two actions
have mutual centralizer algebras.  These statements are
\cite[Theorem~5.18.4, Corollary~5.19.2, and
Theorem~5.22.1]{EtingofEtAl2011}.  They apply over $\C$ in every dimension
used here.  The polynomial modules remain irreducible under the unitary
group: invariance under its differentiated action implies invariance
under its complex span, which is the full matrix Lie algebra.

Let $\pi_{Q,k}^{\lambda}$ denote the Schur-label projector on
$Q^{\otimes k}$.  Define the self-adjoint label observable
\begin{equation}\label{replicas:F-definition}
 F_{Q,k}=\sum_{\lambda\vdash k}(\log d_\lambda)\pi_{Q,k}^{\lambda}.
\end{equation}
We omit the subscript $k$ when the number of copies is unambiguous.
For zero copies, the only label is the empty partition and $F_{Q,0}=0$.
We say labels are compatible if their joint spectral projector is
nonzero.

The next lemma treats three operations on labels.  Merging subsystems
and splitting the copies into groups will be used in the norm comparisons
of Section~\ref{sec:comparators}.  Removing the last copy will control the
ratio between successive metrics below.  The dimension and surprisal
estimates supply entropy bounds for these operations.

\begin{lemma}[Schur-label estimates]\label{lem:schur}
For the finite-dimensional tensor systems just defined, the following
statements hold.
\begin{enumerate}
\item For $\lambda\vdash k$ of length at most $q$, put
$l_i=\lambda_i+q-i$ for $1\le i\le q$.  Then
\begin{equation}\label{replicas:dimensions}
 \dim V^{(q)}_\lambda
   =\prod_{i<j}\frac{l_i-l_j}{j-i},\qquad
 d_\lambda=\frac{k!\prod_{i<j}(l_i-l_j)}{\prod_i l_i!},\qquad
 \log d_\lambda=kH(\lambda/k)+O_q(\log(k+2)).
\end{equation}
The number of labels and the sum of their general-linear multiplicities
are polynomially bounded in $k$, with $q$ fixed.
\item Central label functions commute for nested or disjoint
subsystems.  On $\mathcal S_k$, complementary label projectors agree:
$\pi_{Q,k}^{\lambda}=\pi_{Q^c,k}^{\lambda}$.
\item If $Q,E$ are disjoint and their labels $\lambda,\mu$ are
compatible with label $\nu$ on $QE$, then
\begin{equation}\label{replicas:merge-dimensions}
 d_\nu\le d_\lambda d_\mu,\qquad
 d_\lambda\le d_\nu d_\mu.
\end{equation}
For a tensor product of permutation-invariant density matrices on
$Q^{\otimes k}$ and $E^{\otimes k}$, the conditional probability of
$\nu$ given $\lambda,\mu$ is at most
$\poly(k)d_\nu/(d_\lambda d_\mu)$.  Consequently, for $0\le b\le1$,
\begin{equation}\label{replicas:merge-moment}
 \E\left(\frac{d_\lambda d_\mu}{d_\nu}\right)^b
 \le\poly(k),
\end{equation}
uniformly in those density matrices.
\item Split the copies into specified groups of sizes $k-r$ and $r$.
If their labels $\alpha,\beta$ on a fixed subsystem are compatible with
the full label $\lambda$, then
\begin{equation}\label{replicas:group-dimensions}
 d_\alpha d_\beta\le d_\lambda
       \le \binom{k}{r}d_\alpha d_\beta.
\end{equation}
These subgroup label projectors commute with the full label projector.
\item If $\rho$ is any permutation-invariant density matrix on
$Q^{\otimes k}$ and $L_\rho=-\log\rho$ on its support, then $F_Q$
commutes with $\rho$, $0\le F_Q\le L_\rho$ on that support, and
\begin{equation}\label{replicas:surprisal-moment}
 \Tr\rho\exp\bigl(b(L_\rho-F_Q)\bigr)\le\poly(k)
       \qquad(0\le b\le1).
\end{equation}
\item The full $Q$ label $\lambda$ and the first-$(k-1)$ label $\nu$
are compatible precisely along the possible removable-row branches
$\nu=\lambda-e_i$.  In any permutation-invariant $Q$ state, conditional
on $\lambda$, the probability of this branch is
\begin{equation}\label{replicas:branch-probability}
 p_{\lambda,i}=\frac{d_{\lambda-e_i}}{d_\lambda}
 =\frac{l_i}{k}\frac{\prod_{a<b}(l_a-\delta_{ai}-l_b+\delta_{bi})}
                         {\prod_{a<b}(l_a-l_b)}
 \le 2^q\frac{l_i}{k}.
\end{equation}
The normalized Jucys--Murphy, or star, operator
\cite[Section~2]{VershikOkounkov2005}
\begin{equation}\label{replicas:star-definition}
 J_{Q,k}=\frac1k\sum_{j<k}U_Q((j\ k))
\end{equation}
has eigenvalue $(\lambda_i-i)/k$ on that branch.
\end{enumerate}
The polynomial bounds may depend on the fixed subsystem dimensions,
but not on the density matrices.
\end{lemma}

\begin{proof}
The first formula in \eqref{replicas:dimensions} is the value of the
Weyl character at the identity.  To derive the second, multiply the
character decomposition of $(x_1+\cdots+x_q)^k$ by
$\det[x_i^{q-j}]$ and extract the coefficient of
$x_1^{l_1}\cdots x_q^{l_q}$.  The strict decrease of the $l_i$ selects
only label $\lambda$ and the identity term in its numerator alternant.
On the other side the coefficient is
\[
 k!\det\left[\frac1{(l_i-q+j)!}\right]_{i,j=1}^q
 =\frac{k!}{\prod_i l_i!}
       \det\bigl[(l_i)_{q-j}\bigr]_{i,j=1}^q
 =\frac{k!\prod_{i<j}(l_i-l_j)}{\prod_i l_i!}.
\]
Here a reciprocal negative factorial is zero and $(u)_m$ is the
falling factorial.  The last determinant is a Vandermonde because its
columns are monic polynomials of the indicated degrees.  The estimate
$\log n!=n\log n-n+O(\log(n+2))$ gives the entropy formula, including
zero parts, with $0\log0=0$.  The fixed shifts from $\lambda_i$ to
$l_i$ and the Vandermonde contribute only $O_q(\log(k+2))$.
There are at most $(k+1)^q$ labels, and the dimension formula bounds
each $\dim V^{(q)}_\lambda$ by $(k+q)^{q(q-1)/2}$.

Each label projector is a central character idempotent in the
permutation algebra:
\[
 \pi_{Q,k}^{\lambda}
 =\frac{d_\lambda}{k!}\sum_{\pi\in S_k}
       \chi_\lambda(\pi^{-1})U_Q(\pi).
\]
This follows from character orthogonality, or directly from Schur's
lemma applied to the central sum and its trace on each irreducible.
Conjugation by $U_{QE}(\pi)$ sends a central function of the $Q$
permutations to itself.  This proves nested commutation; disjoint
commutation is immediate.  On $\mathcal S_k$,
$U_Q(\pi)z=U_{Q^c}(\pi^{-1})z$.  A permutation is conjugate to its
inverse, so $\chi_\lambda(\pi)=\chi_\lambda(\pi^{-1})$ and the two
central-projector sums agree there.  This also shows that all Specht
modules are self-dual: their dual characters are the same.

For disjoint subsystems, $[\nu]$ is a constituent of
$[\lambda]\otimes[\mu]$ under diagonal permutations.  This proves
the first inequality of \eqref{replicas:merge-dimensions}.
The identity
\[
 \operatorname{Hom}_{S_k}([\nu],[\lambda]\otimes[\mu])
 \simeq
 \operatorname{Hom}_{S_k}([\lambda],[\nu]\otimes[\mu]^*)
\]
and self-duality prove the second.  Conditional on $\lambda,\mu$, a
product of invariant states is the identity divided by
$d_\lambda d_\mu$ on the two Specht factors.  If
$g_{\lambda\mu\nu}$ is the multiplicity of $[\nu]$ in their tensor
product, its conditional label probability is exactly
$g_{\lambda\mu\nu}d_\nu/(d_\lambda d_\mu)$.
The ambient decomposition of $(Q\otimes E)^{\otimes k}$ implies
\[
 g_{\lambda\mu\nu}\,
 \dim V^{(d_Q)}_\lambda\dim V^{(d_E)}_\mu
 \le \dim V^{(d_Qd_E)}_\nu,
\]
so $g_{\lambda\mu\nu}$ is polynomially bounded.  With
$R=d_\lambda d_\mu/d_\nu\ge1$, the conditional expectation in
\eqref{replicas:merge-moment} is
$\sum_\nu g_{\lambda\mu\nu}R^{b-1}\le\sum_\nu
g_{\lambda\mu\nu}$, again polynomial.

For the copy-group assertion, choose a copy of
$[\alpha]\otimes[\beta]$ inside the restriction of $[\lambda]$ to
$S_{k-r}\times S_r$.  Its dimension gives the lower bound.
The translates of this subspace span $[\lambda]$, by irreducibility.
Choosing one representative of each subgroup coset bounds the span
dimension by $\binom{k}{r}d_\alpha d_\beta$.  Central elements of
$S_k$ commute with its subgroup algebra, proving the commutation claim.

An invariant density has blocks
$\rho=\bigoplus_\lambda I_{[\lambda]}\otimes M_\lambda$.
Every positive eigenvalue $u$ of $M_\lambda$ is repeated $d_\lambda$
times, so $u d_\lambda\le1$.  This gives $F_Q\le L_\rho$ and
\[
 \Tr\rho e^{b(L_\rho-F_Q)}
 =\sum_{\lambda}\sum_{u>0}(u d_\lambda)^{1-b}
 \le\sum_\lambda\dim V^{(d_Q)}_\lambda
 \qquad(0\le b\le1),
\]
where the inner sum counts eigenvalues with multiplicity in
$M_\lambda$.

It remains to identify the marked-copy branches.  Multiplying a Weyl
numerator alternant by $\sum_i x_i$ gives the sum of alternants with
one exponent increased by one.  Terms with equal exponents vanish;
the remaining terms correspond exactly to adding a legal box to a
partition.  Thus
$V^{(q)}_\nu\otimes\C^q=\bigoplus_{\lambda=\nu+e_i}
V^{(q)}_\lambda$, with multiplicity one.  Comparing the two
Schur--Weyl decompositions of the first $k-1$ copies and the last copy
shows that the $\nu$ branch in the $\lambda$ block has rank
$d_\nu$ on $[\lambda]$, and acts identically on $V^{(q)}_\lambda$.
The invariant state is uniform on $[\lambda]$, giving the probability
in \eqref{replicas:branch-probability}.  The dimension formula gives
the displayed quotient.  In its Vandermonde ratio, a factor involving
an earlier row is at most two, and a factor involving a later row is
at most one.  Positivity holds on a legal removable branch.

For the star eigenvalue, let $E_{ab}^{\rm tot}$ be the sum over copies
of the matrix unit $E_{ab}$ on $\C^q$.  Direct multiplication gives
\[
 \sum_{a,b}E_{ab}^{\rm tot}E_{ba}^{\rm tot}
 =kq I+2\sum_{j<j'}U_Q((j\ j')).
\]
The right side is scalar on a Schur block.  The lexicographically
leading weight in its Weyl character is $\lambda$; a vector of that
weight is annihilated by $E_{ab}^{\rm tot}$ for $a<b$, since those
operators would increase the weight.  On this vector the diagonal
terms contribute $\sum_i\lambda_i^2$, and the two terms for $a<b$
contribute $\lambda_a-\lambda_b$, using
$[E_{ab}^{\rm tot},E_{ba}^{\rm tot}]=E_{aa}^{\rm tot}-E_{bb}^{\rm tot}$.
The scalar is therefore
$\sum_i\lambda_i(\lambda_i+q+1-2i)$.
Subtract the formula for $k-1$ copies with label $\lambda-e_i$.
The difference of the transposition sums is $kJ_{Q,k}$, and its value
is $\lambda_i-i$.
\end{proof}

\subsection{A common metric on every subsystem}

Fix $0<t<1/4$ and set $a=2t$.
Complementary labels will need to give exactly the same metric, even
when the two subsystem dimensions differ.  We therefore use the full
dimension $d$ in the following formula for every subsystem.  Besides
having eigenvalues of size $e^{-tF_Q}$ up to polynomial factors, the
metric must admit a positive mixture of tensor powers, so that the
conditional movement estimate can be applied inside its integral.
The ratio between its $k$-copy and $(k-1)$-copy weights must also
approximate marginal powers for the local-operator calculation.
The following common label function has all these properties.
Pad each label to length $d$ and put
\begin{equation}\label{replicas:w-definition}
 \begin{split}
 c&=d+(1+t)\binom{d}{2},\\
 w_k(\lambda)&=\frac{\Gamma(c)}{\Gamma(c+tk)}
    \prod_{i=1}^d
       \frac{\Gamma(1+t(\lambda_i+d-i))}
            {\Gamma(1+t(d-i))},\\
 W_{Q,k}&=\sum_\lambda w_k(\lambda)\pi_{Q,k}^{\lambda}.
 \end{split}
\end{equation}
In particular $W_{Q,0}=I$.  For a one-dimensional subsystem at
positive $k$, $W_{Q,k}$ is the positive scalar $w_k((k))$; it need
not equal one.  This convention keeps the label function identical
on all subsystems.

\begin{lemma}[Integral formula and marked ratios]\label{replicas:metric}
The operators in \eqref{replicas:w-definition} are positive definite
and central.  Uniformly in their labels,
\begin{equation}\label{replicas:W-comparison}
 \poly(k)^{-1}e^{-tF_{Q,k}}\le W_{Q,k}
       \le\poly(k)e^{-tF_{Q,k}}.
\end{equation}
On $\mathcal S_k$, $W_{Q,k}=W_{Q^c,k}$ exactly.  For each subsystem
there is a probability measure on positive definite subdensities
$\sigma$ on $Q$, with $\Tr\sigma\le1$, such that
\begin{equation}\label{replicas:W-integral}
 W_{Q,k}=\int(\sigma^t)^{\otimes k}\,d\zeta_Q(\sigma)
                  \qquad(k\ge0).
\end{equation}
The measure is independent of $k$.

Extend $W_{Q,k-1}$ by the identity on the last copy and define
$R_{Q,k}=W_{Q,k}W_{Q,k-1}^{-1}$.  Then
\begin{equation}\label{replicas:R-forward}
 \norm{R_{Q,k}-(J_{Q,k})_+^t}\longrightarrow0
\end{equation}
on the entire tensor space.  There is a constant $C_{\mathcal V,t}$
such that, for every unit $z\in\mathcal S_k$,
$\norm{R_{Q,k}^{-1}z}\le C_{\mathcal V,t}$, and
\begin{equation}\label{replicas:R-inverse}
 \limsup_{k\to\infty}\sup_{\substack{z\in\mathcal S_k\\\norm z=1}}
 \norm{\bigl(R_{Q,k}^{-1}-\max(J_{Q,k},\delta)^{-t}\bigr)z}
 \le C_{\mathcal V,t}\delta^{(1-2t)/2}
                 \qquad(0<\delta<1).
\end{equation}
\end{lemma}

\begin{proof}
Begin in dimension $d$.  Write $\Delta(s)=\prod_{i<j}(s_i-s_j)$
and choose $s$ on the simplex $s_i\ge0$, $\sum_i s_i=1$, with density
proportional to $\Delta(s)\Delta(s^t)$.  This is a nonnegative,
integrable density, positive off the collision hyperplanes in the
simplex interior.  It is the simplex normalization of the Laguerre
Muttalib--Borodin weight; the gamma integral below is the corresponding
form of \cite[Section~2.3, equation~(26)]{ForresterIpsen2018}.
We include the calculation before verifying the common-subsystem and
ratio properties.  For $d=1$ use the point mass at $s_1=1$.
Choose an independent Haar-distributed unitary eigenbasis and let
$\sigma$ have eigenvalues $s$.  The integral of
$(\sigma^t)^{\otimes k}$ commutes both with permutations and with
the diagonal unitary group; hence it is central in the Schur
decomposition.  Its eigenvalue on $\lambda$ is the mean of
$\chi_\lambda^{(d)}(s^t)/\dim V^{(d)}_\lambda$.

Put $l_i=\lambda_i+d-i$.  The Weyl denominator cancels
$\Delta(s^t)$ in the density.  Expanding the two remaining
determinants and using the simplex integral
\[
 \int_{\substack{s_i\ge0\\\sum s_i=1}}
        \prod_i s_i^{u_i}\,ds
     =\frac{\prod_i\Gamma(1+u_i)}{\Gamma(d+\sum_i u_i)}
                 \qquad(u_i>-1)
\]
gives, before dividing by the Weyl dimension and the density's
normalization,
\[
 \frac{d!}{\Gamma(c+tk)}
       \det\bigl[\Gamma(1+t l_i+d-j)\bigr]_{i,j=1}^d.
\]
The simplex identity follows, for example, by changing variables
$x_i=rs_i$ in the product of the defining gamma integrals; the
Jacobian contributes $r^{d-1}$.  Factor $\Gamma(1+t l_i)$ from
each row.  The remaining entries are rising-factorial monic
polynomials in $t l_i$ of degrees $d-j$, so their determinant is
$t^{d(d-1)/2}\Delta(l)$.  The Weyl dimension cancels $\Delta(l)$.
Normalizing at $k=0$ gives exactly \eqref{replicas:w-definition}.

For $d_Q\le d$, choose an isometry $E:\mathcal H_Q\to\C^d$.
Compression of a permutation on $(\C^d)^{\otimes k}$ by
$E^{\otimes k}$ is the same permutation on $Q^{\otimes k}$.
Compressing its central label function therefore gives the identical
function $w_k(\lambda)$ on all labels that occur in dimension $d_Q$.
On each sample, write $B=E^\dagger\sigma^t E$ and
$\tau=B^{1/t}$.  The min--max principle bounds the decreasing
eigenvalues of $B$ by the corresponding eigenvalues of $\sigma^t$.
Consequently $\Tr\tau\le\Tr\sigma=1$, while $B>0$ almost surely.
Thus the compressed integral has the form
\eqref{replicas:W-integral}, with the distribution of $\tau$ as
$\zeta_Q$.  This proves the asserted representation without changing
the common label function.  Exact complement equality follows from
Lemma~\ref{lem:schur}.

The gamma estimate
$\log\Gamma(1+u)=u\log u-u+O(\log(u+2))$, with $0\log0=0$,
applied to the fixed number of factors in
\eqref{replicas:w-definition}, gives
$\log w_k(\lambda)=-tkH(\lambda/k)+O_{d,t}(\log(k+2))$.
The terms containing $\log t$ cancel.  Formula
\eqref{replicas:dimensions} proves \eqref{replicas:W-comparison}.

For the marked ratio, now put $l_i=\lambda_i+d-i$ using the full
dimension.  On a removable branch its exact eigenvalue is
\begin{equation}\label{replicas:ratio-eigenvalue}
 r_{\lambda,i}=
 \frac{\Gamma(c+t(k-1))}{\Gamma(c+tk)}
 \frac{\Gamma(1+t l_i)}{\Gamma(1+t(l_i-1))}.
\end{equation}
Gamma ratios give $r_{\lambda,i}\asymp_{d,t}((l_i+1)/k)^t$.
When $l_i/k$ is bounded away from zero, the same ratios give uniform
convergence to $(l_i/k)^t$.  When $l_i/k$ is small, both
$r_{\lambda,i}$ and $\max((l_i-d)/k,0)^t$ are uniformly small.
First separate at $l_i/k=\varepsilon$, then let
$k\to\infty$ and $\varepsilon\downarrow0$.  The star eigenvalue
formula proves \eqref{replicas:R-forward} and a uniform bound for
$\norm{R_{Q,k}}$.

The $Q$ marginal of a symmetric vector is permutation invariant.
Its conditional branch probabilities are bounded by $C_d l_i/k$:
Lemma~\ref{lem:schur} first uses $\lambda_i+d_Q-i$, which is no
larger than $l_i$.  Hence, conditional on any $\lambda$,
\[
 \sum_i p_{\lambda,i}r_{\lambda,i}^{-2}
 \le C_{d,t}\sum_i(l_i/k)^{1-2t}\le C_{d,t}.
\]
This proves the inverse vector bound.  On the spectral part
$J_{Q,k}\ge\delta$, the inverse approximation follows in operator
norm from \eqref{replicas:R-forward}.  On the remaining part,
$l_i/k\le\delta+d/k$.  The true inverse's squared norm there is at
most $C_{d,t}(\delta+d/k)^{1-2t}$; the clipped inverse's squared
norm is at most $C_d(\delta+d/k)\delta^{-2t}$.
Taking the limit superior and the square root proves
\eqref{replicas:R-inverse}.  All bounds were conditional on an
arbitrary label, so they are uniform in the original symmetric
vector.
\end{proof}

The common label function also gives an exact restriction identity.
If $B,Q$ are disjoint and $B$ is held in a fixed pure vector
$\phi^{\otimes k}$, then
\begin{equation}\label{replicas:fixed-factor}
 W_{BQ,k}(\phi^{\otimes k}\otimes z)
       =\phi^{\otimes k}\otimes W_{Q,k}z.
\end{equation}
Indeed every simultaneous permutation on $BQ$ fixes the repeated
$B$ vector and acts as the corresponding permutation on $Q$.
Apply the same central permutation coefficients on both sides.
The corresponding identity holds for $F$ and for every real power of
$W$.  In particular it applies when the fixed vector on $B$ is a Bell
vector between two factors: its entanglement inside $B$ is irrelevant
to simultaneous permutations of the entire $B$ factor.

For a partition $P,Y,F$ of the full one-copy system, the moves considered
below transfer a subsystem $x\subseteq Y$ to either $P$ or $F$.
We call $P,F$ the outer parts and $Y$ the middle part; these names
impose no geometric condition. Define the partition's metric on
$\mathcal S_k$ by
\begin{equation}\label{replicas:leaf-metric}
 A_k(P,Y,F)=\bigl(W_{P,k}^{-1}W_{F,k}^{-1}W_{Y,k}\bigr)^2.
\end{equation}
The three factors commute.  The preceding lemmas imply
\begin{equation}\label{replicas:leaf-floor}
 \log A_k=a(F_P+F_F-F_Y)+O_{\mathcal V,t}(\log(k+2)),
 \qquad A_k\ge\poly(k)^{-1}I.
\end{equation}
The error here is an operator-norm bound.  Indeed $F_Y=F_{PF}$ on
symmetry and $F_{PF}\le F_P+F_F$ on every compatible label.
If a one-copy operator $h$ is contained in one part, its copy average
$\bar h=k^{-1}\sum_jh^{(j)}$ commutes with this metric: it commutes
with permutations on that part and is disjoint from the other parts.

\subsection{An operator lower bound for moving a subsystem}

The integral representation of $W$ lets us apply the one-copy movement
estimate inside a replica metric.  We first show that moving $x$ from
the middle part to an outer part produces the corresponding entropy gain
as a rank-one operator bound.  For a unit one-copy vector $\theta$, put
$P_{\theta,k}=\ket{\theta^{\otimes k}}\bra{\theta^{\otimes k}}$.
All constants outside polynomial factors in $k$ will depend only on
the one-copy estimate and the dimension of $x$.

\begin{lemma}[Relative coherent pin]\label{lem:relative-pin}
Let $P,x,Y_0,F$ partition the fixed finite system $\mathcal V$, put
$Y=xY_0$, and let
\[
 A_k=A_k(P,xY_0,F),\qquad B_k=A_k(Px,Y_0,F),\qquad
 C_{x,k}=A_k^{-1/2}B_kA_k^{-1/2}.
\]
Suppose $0<t<1/4$, $a=2t$, and
$a\log(e\dim x)\le c$, with $c$ as in Lemma~\ref{lem:movement}.
For every unit $\theta\in\mathcal V$, define
\[
 \eta_\theta=S_\theta(x|P)+S_\theta(x|Y_0)
             =I_\theta(x:F|P)=I_\theta(x:F|Y_0)\ge0.
\]
There is a polynomial bound uniform in $\theta$ such that, on
$\mathcal S_k$,
\begin{equation}\label{replicas:relative-pin}
 C_{x,k}\ge\poly(k)^{-1}
   \exp\left\{ka\left[\eta_\theta
          -C a^{1/4}\log^C(e\dim x)\right]\right\}
             P_{\theta,k}.
\end{equation}
The same statement holds for a move to $F$ after exchanging the outer
parts.  A choice making no move has $C_{x,k}=I$ and may be assigned
gain zero.
\end{lemma}

\begin{proof}
For a positive definite $C$ and unit vector $z$, Cauchy--Schwarz gives
\[
 |\langle z,u\rangle|^2
 \le\langle z,C^{-1}z\rangle\langle u,Cu\rangle,
 \qquad
 C\ge\langle z,C^{-1}z\rangle^{-1}\ket{z}\bra{z}.
\]
It therefore suffices to bound the inverse quadratic form of
$C_{x,k}$ at $\theta^{\otimes k}$.  Since
$C_{x,k}^{-1}=A_k^{1/2}B_k^{-1}A_k^{1/2}$,
\begin{equation}\label{replicas:move-inverse}
 \begin{split}
 \langle\theta^{\otimes k},C_{x,k}^{-1}\theta^{\otimes k}\rangle
 &=\norm{B_k^{-1/2}A_k^{1/2}\theta^{\otimes k}}^2\\
 &=\norm{W_{Px,k}W_{xY_0,k}W_{P,k}^{-1}W_{Y_0,k}^{-1}
                                      \theta^{\otimes k}}^2.
 \end{split}
\end{equation}
To obtain the second line, cancel the disjoint $F$ factors and commute
the remaining nested or disjoint central factors.  This is an
identity of the displayed norms; no commutation of $A_k$ with $B_k$
is asserted.

Let $\rho_R$ be the marginal of $\theta$ for $R=P,Y_0$, and set
\[
 T_R=W_{R,k}^{-1}(\rho_R^t)^{\otimes k},
\]
with the positive power zero on the kernel.  The invariant density
$\rho_R^{\otimes k}$ has every eigenvalue in its $\lambda$ block
at most $d_\lambda^{-1}$ by eigenvalue repetition, as in the proof
of \eqref{replicas:surprisal-moment}.  Thus
\eqref{replicas:W-comparison} implies
\begin{equation}\label{replicas:compensator}
 \norm{T_R}\le\poly(k),
\end{equation}
uniformly in $\theta$, including marginals of deficient rank.
Furthermore $T_P$ commutes with both numerator factors in
\eqref{replicas:move-inverse}, and so does $T_{Y_0}$.
For the overlapping numerator this follows because $W_{R,k}$ is
central on the nested subsystem and $(\rho_R^t)^{\otimes k}$
commutes with every simultaneous permutation on $xR$; the other
numerator is disjoint from $R$.

Extend $\rho_R$ by eigenvalue one on its kernel and denote this
extension by $\widetilde\rho_R$.  The product of the two marginal
support projectors fixes $\theta$, since $P$ and $Y_0$ are disjoint.
We can therefore insert $T_P,T_{Y_0}$ and move them to the left in
\eqref{replicas:move-inverse}, obtaining the upper bound
\[
 \poly(k)\,
 \norm{W_{Px,k}W_{xY_0,k}
       (\widetilde\rho_P^{-t}\widetilde\rho_{Y_0}^{-t})^{\otimes k}
                                       \theta^{\otimes k}}^2.
\]
Here and below the polynomial can be enlarged to absorb its square.
Expand the two numerator integrals
\eqref{replicas:W-integral} in their stated order.  The triangle
inequality and the tensor-power norm identity bound this by
\begin{equation}\label{replicas:move-one-copy}
 \poly(k)\sup_{\sigma,\tau}
 \norm{\sigma_{Px}^t\tau_{xY_0}^t
       \widetilde\rho_P^{-t}\widetilde\rho_{Y_0}^{-t}\theta}^{2k},
\end{equation}
where $\sigma,\tau$ are positive definite subdensities on the
indicated subsystems.  Normalizing a subdensity multiplies its
$t$-power by a scalar at least one, so it suffices to take density
matrices.  Since $P$ is disjoint from $xY_0$, its inverse power
commutes with $\tau_{xY_0}^t$.  The one-copy word is therefore
\[
 \sigma_{Px}^{a/2}\widetilde\rho_P^{-a/2}
 \tau_{xY_0}^{a/2}\widetilde\rho_{Y_0}^{-a/2}\theta,
\]
exactly the word in Lemma~\ref{lem:movement}.
That lemma bounds its log norm by
$-a\eta_\theta/2+C a^{5/4}\log^C(e\dim x)$.
Substituting in \eqref{replicas:move-one-copy}, and then using the
inverse Cauchy--Schwarz inequality at the start of the proof, gives
\eqref{replicas:relative-pin}.  Every polynomial bound used in this
argument was uniform in $\theta$.
\end{proof}

\subsection{Scalar symbols, products, and absolute values}

The relative lower bound is now available.  To control energy, we also
need scalar descriptions of products and absolute values of local terms
after a metric similarity transform.  This requires a further argument
beyond their individual expectations on product vectors.

Let $d\theta$ be invariant probability measure on the unit sphere of
$\mathcal V$.  Unitary invariance and irreducibility of $\mathcal S_k$
give
\begin{equation}\label{replicas:resolution}
 \int P_{\theta,k}\,d\theta=D_k^{-1}\Pi_k.
\end{equation}
The trace fixes the scalar.  For a density matrix $\sigma$ supported
on $\mathcal S_k$, define its coherent probability measure by
\begin{equation}\label{replicas:husimi}
d\mu_\sigma(\theta)=D_k\Tr(\sigma P_{\theta,k})\,d\theta.
\end{equation}
This is the coherent measure used in finite quantum de~Finetti
theorems~\cite{ChristandlEtAl2007,LewinNamRougerie2015}.  We reproduce its
fixed-copy approximation in the proof below, then establish the additional
product calculus needed for the nonlinear functions of our metric.
Here $\theta$ is a pure one-copy state, whereas $\sigma$ may be any
symmetric replica density.  The same $\sigma$ defines both the operator
expectation and the measure mixing the one-copy symbols below.

\begin{lemma}[Marked-copy scalar symbols]\label{lem:symbol}
Fix a partition $P,Y,F$ of $\mathcal V$, $0<t<1/4$, and a one-copy
operator $h$ supported on $PY$.  Let $A_k$ be
\eqref{replicas:leaf-metric}, and define on $\mathcal S_k$
\[
 O_k=A_k^{-1/2}\bar h A_k^{1/2},\qquad
 f_\theta=\langle\theta,
       \rho_P^t\rho_Y^{-t}h\rho_P^{-t}\rho_Y^t\theta\rangle,
\]
where $\rho_Q$ is the $Q$ marginal of $\theta$.  All powers in
$f_\theta$ are zero on kernels.  Then $\sup_k\norm{O_k}<\infty$,
with a bound allowed to depend on the fixed entire system and $h$.
For every fixed noncommutative polynomial $q$ in two variables,
\begin{equation}\label{replicas:polynomial-symbol}
 \sup_{\sigma}
 \left|\Tr\sigma q(O_k,O_k^\dagger)
       -\int q(f_\theta,\overline{f_\theta})\,d\mu_\sigma(\theta)
 \right|\longrightarrow0,
\end{equation}
where the supremum is over all density matrices on $\mathcal S_k$.
Moreover, for $\mathcal D_k=|O_k|+|O_k^\dagger|-O_k-O_k^\dagger$,
\begin{equation}\label{replicas:skew-symbol}
 \Tr\sigma\mathcal D_k
 =2\int\bigl(|f_\theta|-\operatorname{Re}f_\theta\bigr)
                       \,d\mu_\sigma(\theta)+o_k(1),
\end{equation}
uniformly in $\sigma$.

If $0\le h\le I$ has a designated supporting subsystem of dimension
$d_h$, put $x=Y\cap\supp h$, $Y_0=Y\setminus x$, and
\[
 \eta_\theta=S_\theta(x|P)+S_\theta(x|Y_0)
            =I_\theta(x:F|P)=I_\theta(x:F|Y_0).
\]
Here the designated support may replace the minimal support in the
definition of $x$.  If $a\log(e d_h)\le c$, then
\begin{equation}\label{replicas:skew-bound}
 \Tr\sigma\mathcal D_k
 \le C a^2\log^C(e d_h)
                  \int\eta_\theta^{1/8}\,d\mu_\sigma(\theta)
                   +o_k(1).
\end{equation}
The assertion also holds with the two outer parts exchanged.
\end{lemma}

\begin{proof}
We first replace the marked ratios by bounded polynomial words, then
prove a product calculus for those words, and finally pass to absolute
values.  The inverse approximation must be applied on symmetric
input vectors; no full-space bound for $R_{Q,k}^{-1}$ is needed.

\paragraph{Bounded approximation on symmetric inputs.}
Between symmetric vectors, $h^{(k)}$ may replace $\bar h$.
The $F$ factors commute with this term and cancel.  Since
$W_{Q,k-1}$ commutes with $W_{Q,k}$ and is disjoint from $h^{(k)}$,
the remaining form is the compression of
\begin{equation}\label{replicas:marked-word}
 R_{P,k}R_{Y,k}^{-1}h^{(k)}R_{P,k}^{-1}R_{Y,k}.
\end{equation}
Write $h$ as a finite sum of products $c_P\otimes e_Y$.
Such a decomposition follows by expanding matrix units on the two
parts of its support; its length and norm bounds are fixed here.
For one product and a symmetric unit vector $z$, the quadratic form
of \eqref{replicas:marked-word} equals
\begin{equation}\label{replicas:inverse-innermost}
 \left\langle R_{Y,k}e_Y^\dagger R_{Y,k}^{-1}z,
                  R_{P,k}c_P R_{P,k}^{-1}z\right\rangle.
\end{equation}
Both inverses act directly on $z$.  The uniform forward operator bound
and inverse vector bound from Lemma~\ref{replicas:metric} bound these
two vectors uniformly and therefore bound the numerical radius of
$O_k$.  The elementary inequality
$\norm T\le2\sup_{\norm z=1}|\langle z,Tz\rangle|$, obtained by
polarizing the quadratic form, proves uniform boundedness of $O_k$.

In \eqref{replicas:inverse-innermost}, replace each inverse by
$\max(J_{Q,k},\delta)^{-t}$ using
\eqref{replicas:R-inverse}; then replace each forward ratio by
$(J_{Q,k})_+^t$ using \eqref{replicas:R-forward}.
The first replacement has limiting error at most
$C_{\mathcal V,t,h}\delta^{(1-2t)/2}$, uniformly on symmetric unit
vectors.  At fixed $\delta$ the second has error tending to zero.
Approximate the continuous functions $u\mapsto u_+^t$ and
$u\mapsto\max(u,\delta)^{-t}$ uniformly on $[-1,1]$ by polynomials.
All star spectra lie in this interval.  Polarization, applied after
summing the finitely many product components, gives compressed
operators $B_k$, each a fixed polynomial word in marked operators
and stars, such that
\begin{equation}\label{replicas:polynomial-approximation}
 \limsup_k\norm{O_k-B_k}\le\varepsilon.
\end{equation}
Here any $\varepsilon>0$ is achieved by first choosing $\delta$,
then the polynomials, and then taking the limit in $k$.
The approximants may be chosen with a common bound on
$\limsup_k\norm{B_k}$; an approximation-dependent threshold for $k$
does not affect any of the limits below.

\paragraph{The coherent measure on a fixed number of copies.}
We next establish the uniform coherent-measure calculation for a
fixed number of copies.  If $G$ acts on $\mathcal V^{\otimes m}$,
let $\mathcal T_{k,m}(G)$ be its average over all ordered injections
of these $m$ copies into $k$ copies.  It preserves $\mathcal S_k$.
For $k\ge m$, symmetry gives
$\Tr\sigma\mathcal T_{k,m}(G)=\Tr\sigma^{(m)}G$.
For the integral of its product-vector expectation,
\eqref{replicas:resolution} gives the finite-copy identity
\cite[equation~(5)]{Chiribella2011}
\begin{equation}\label{replicas:finite-husimi}
 \int\langle\theta^{\otimes m},G\theta^{\otimes m}\rangle
                         \,d\mu_\sigma(\theta)
 =\frac{D_k}{D_{k+m}}
       \Tr\bigl[(\sigma\otimes G)\Pi_{k+m}\bigr].
\end{equation}
In the permutation expansion of $\Pi_{k+m}$, the fraction mapping
every one of the last $m$ indices into the first $k$ is
$(k)_m/(k+m)_m=1-O_m(k^{-1})$.  Any such permutation is equivalent,
by a permutation of the first $k$ indices on each side, to $m$
disjoint swaps matching last indices with distinct first indices.
Indeed reorder the old images of the last indices on the left and
their old preimages on the right, then eliminate the remaining
permutation among old indices.  Both old permutations are absorbed
because $U(\pi)\sigma=\sigma U(\pi)=\sigma$.
The matched swaps contribute exactly $\Tr\sigma^{(m)}G$ by partial
trace.  Every other term has absolute value at most
$\norm{\sigma\otimes G}_1\le d^m\norm G$.
Also $D_k/D_{k+m}=1+O_{d,m}(k^{-1})$.  Consequently
\begin{equation}\label{replicas:uniform-husimi}
 \left|\Tr\sigma\mathcal T_{k,m}(G)
  -\int\langle\theta^{\otimes m},G\theta^{\otimes m}\rangle
                                      \,d\mu_\sigma\right|
 \le\frac{C_{d,m}\norm G}{k},
\end{equation}
uniformly in $\sigma$.

\paragraph{Products after compression.}
The remaining issue is multiplication after compression.  For any
ambient operator $B$, define its full-copy permutation average
\[
 \mathcal T_k(B)=\frac1{k!}\sum_{\pi\in S_k}U(\pi)BU(\pi)^\dagger.
\]
It commutes with all $U(\pi)$, hence preserves symmetry, and
\[
 \Pi_k B\Pi_k\big|_{\mathcal S_k}
   =\mathcal T_k(B)\big|_{\mathcal S_k}.
\]
Products of compressed words are therefore restrictions of products
of their full permutation averages.
The multiplication of these averages uses disjoint placements of fixed
copy blocks, as in the mean-field symmetrization argument of
\cite[Lemma~IV.1]{RaggioWerner1989}.  For fixed operators $G,H$ on
$m,n$ copies, respectively,
\begin{equation}\label{replicas:injection-product}
 \left\|\mathcal T_{k,m}(G)\mathcal T_{k,n}(H)
       -\mathcal T_{k,m+n}(G\otimes H)\right\|
 \le \frac{C_{m,n}\|G\|\|H\|}{k}
\end{equation}
for sufficiently large $k$.  Indeed two independently chosen injections
have intersecting images with probability $O_{m,n}(k^{-1})$.
Conditional on disjoint images, their union is a uniform ordered
injection of $m+n$ copies.  The intersecting terms and the change in
normalization each cost at most that probability times $\|G\|\|H\|$.

Expand a fixed star word as sums over its donor indices, with one
marked center.  Repeated donor indices have total coefficient weight
$O(k^{-1})$, so their contribution has norm $O(k^{-1})$ for that
fixed word.  After averaging each word separately, both its center
and its donors occupy a uniformly relocated block of indices.
For a monomial containing $r$ stars, its distinct-donor terms give
an injection average with coefficient $(k-1)_r/k^r=1+O_r(k^{-1})$.
Thus each averaged word is, up to operator-norm error $O(k^{-1})$,
a finite sum of injection averages of fixed operators.
Iterating \eqref{replicas:injection-product} puts a product of words
on disjoint blocks, with a separate center for each word.

On a product vector, a donor appearing in one partial swap contracts
to $\rho_Q$ on the center.  In matrix units this is the identity
\[
 \Tr_{\rm donor}\left[
  (\ket{\theta}\bra{\theta})_{\rm donor}
  \sum_{a,b}(E_{ab})_{Q,\rm donor}\otimes(E_{ba})_{Q,\rm center}
                      \right]=(\rho_Q)_{\rm center}.
\]
Tracing distinct donors successively preserves the order of the
operators on each center.  The separate centers give multiplication
of scalar expectations, since
\[
 \langle\theta^{\otimes(m+n)},(G\otimes H)
                         \theta^{\otimes(m+n)}\rangle
 =\langle\theta^{\otimes m},G\theta^{\otimes m}\rangle
  \langle\theta^{\otimes n},H\theta^{\otimes n}\rangle.
\]
It follows that a polynomial approximant has the
symbol obtained by replacing its stars by the corresponding marginal
densities, its adjoint has the complex conjugate symbol, and products
multiply these scalar symbols.  Applying
\eqref{replicas:uniform-husimi} proves these assertions uniformly in
the symmetric density matrix, for each fixed polynomial word.

\paragraph{Removing clipping and taking absolute values.}
We have proved the product assertion with fixed clipping.  To remove
clipping on the one-copy side, use
\[
 \sum_{0<p_i<\delta}p_i^{1-2t}
       \le d_Q\delta^{1-2t}
\]
for the eigenvalues of any marginal density.  This bounds the squared
norm of its inverse-power tail on $\theta$, uniformly in $\theta$.
The clipped one-copy vectors corresponding to
\eqref{replicas:inverse-innermost} therefore converge uniformly to their
support-power limits, and their inner product converges uniformly to
the zero-kernel-power expression $f_\theta$.
In particular this
symbol is bounded and continuous, as a uniform limit of the clipped
continuous symbols.  Telescoping fixed products in
\eqref{replicas:polynomial-approximation}, using the common limiting
norm bound, now proves \eqref{replicas:polynomial-symbol}.

In particular, both $O_k^\dagger O_k$ and $O_kO_k^\dagger$ have
symbol $|f_\theta|^2$.  Choose an interval $[0,M^2]$ containing all
their spectra and all values $|f_\theta|^2$.  Uniform polynomial
approximation of the square root on that interval gives symbols
$|f_\theta|$ for $|O_k|$ and $|O_k^\dagger|$, with arbitrarily small
uniform expectation errors.  This argument applies whether or not
$O_k$ is normal.  Combining with the linear terms proves
\eqref{replicas:skew-symbol}.  Finally apply
Lemma~\ref{lem:skew} pointwise to $f_\theta$ to obtain
\eqref{replicas:skew-bound}.  Purity gives the two expressions for
$\eta_\theta$; enlarging $x$ to the designated supporting subsystem
is allowed by that lemma.  The proof is unchanged after exchanging
$P$ and $F$.
\end{proof}

The section has supplied two compatible one-copy quantities: the
gain $\eta_\theta$ in \eqref{replicas:relative-pin}, and the cost
$\eta_\theta^{1/8}$ in \eqref{replicas:skew-bound}.  The operator lower
bound can be tested against any symmetric density, and the symbol
error is uniform over those densities, with the prescribed order of
limits.  They can therefore be applied to the same density matrices
when metric changes are averaged.
